\begingroup
\shorthandoff{"<>"}
\setlength{\arrayrulewidth}{1pt}
\arrayrulecolor{vlepsbased}
\begin{longtable}{|p{15cm}|}
  \hline
  \rowcolor{maincolor} 
  \multicolumn{1}{|l|}{\textcolor{white}{\textbf{Caso de Uso}}} \\
  \hline
  \endfirsthead
  \multicolumn{1}{c}{{\tablename\ \thetable{} -- Continuación}} \\
  \hline
  \rowcolor{maincolor} 
  \multicolumn{1}{|l|}{\textcolor{white}{\textbf{Caso de Uso (Continuación)}}} \\
  \hline
  \endhead
  \hline \multicolumn{1}{|r|}{\textit{Continúa en la siguiente página...}} \\ \hline
  \endfoot
  \hline
  \caption{Caso de Uso CU5 — Inicio de sesión en la plataforma.}
  \label{TB:CU5_LOGIN}
  \endlastfoot

  \rowcolor{ulepsbased}
  \begin{tabularx}{\linewidth}{@{}p{7.5cm}p{7.5cm}@{}}
    \textbf{Identificador:} CU5 & \textbf{Actor principal:} Usuario
  \end{tabularx} \\ \hline

  \textbf{Nombre:} Inicio de sesión en la plataforma \\ \hline

  \rowcolor{ulepsbased}
  \textbf{Descripción:} \newline 
  El usuario se autentica en la plataforma mediante sus credenciales para acceder al entorno de tutoría. \\ \hline

  \textbf{Precondiciones:} \newline
  1. El usuario debe existir en el sistema. \\ \hline

  \rowcolor{ulepsbased}
  \textbf{Flujo principal:} \newline
  1. El usuario accede a la vista de "Inicio de sesión". \newline
  2. Introduce su correo electrónico o usuario y su contraseña. \newline
  3. Pulsa "Ingresar". \newline
  4. El sistema valida las credenciales verificando el hash almacenado. \newline
  5. El sistema genera los tokens JWT de acceso y refresco correspondientes. \newline
  6. El sistema devuelve los tokens y redirige al panel principal del tutor. \\ \hline

  \textbf{Flujos alternativos:} \newline
  \textbf{FA-01:} Credenciales incorrectas: \newline
  1. El sistema detecta que el usuario no existe o la contraseña no es válida. \newline
  2. Se deniega el acceso, se no generan tokens y se muestra un error al usuario. \\ \hline

  \rowcolor{ulepsbased}
  \textbf{Postcondiciones:} \newline
  1. Se establece un contexto de sesión segura. \newline
  2. El usuario tiene acceso completo a las áreas privadas. \\ \hline

\end{longtable}
\endgroup
